% TOP_PROBLEM: {"problemId": "problem.top500-omission-w045049-exponential-algebraic-closedness", "canonicalTitle": "Zilber's Exponential-Algebraic Closedness Conjecture", "releaseRank": 441, "primaryDomain": "logic_foundations_set_theory", "primaryDomainLabel": "Logic, foundations & set theory", "displayStatus": "open_with_solved_subcases", "releaseStatus": "open_with_solved_subcases"}
% !TeX program = lualatex
% Definition and sources only; this file is not an LLM solution attempt.
\documentclass[11pt]{article}
\usepackage[margin=25mm]{geometry}
\usepackage{fontspec}
\setmainfont{Latin Modern Roman}
\usepackage{amsmath,amssymb,mathtools}
\usepackage{unicode-math}
\setmathfont{Latin Modern Math}
\usepackage{xurl,hyperref}
\hypersetup{colorlinks=true,urlcolor=blue}
\setcounter{secnumdepth}{0}
\setlength{\parindent}{0pt}
\setlength{\parskip}{5pt}
\setlength{\emergencystretch}{3em}
\begin{document}
\section*{\texorpdfstring{Zilber's Exponential-Algebraic Closedness Conjecture}{Zilber's Exponential-Algebraic Closedness Conjecture}}\label{441}
\subsection{Definitions and mathematical statement}
Work in ${\mathbb{C}}^n\times({\mathbb{C}}^*)^n$. A variety $V$ is free if its additive projection is not contained in a translate of a proper rational linear subspace and its multiplicative projection is not contained in a translate of a proper algebraic subtorus. For a linear subspace $L\subseteq{\mathbb{C}}^n$ defined over ${\mathbb{Q}}$, the coordinatewise exponential image $\exp(L)$ is an algebraic subtorus. Write $\pi_L$ for the quotient map to
\[
 ({\mathbb{C}}^n/L)\times(({\mathbb{C}}^*)^n/\exp(L)).
\]
Call $V$ rotund if $\dim\pi_L(V)\ge n-\dim L$ for every such $L$. The conjecture asks that every irreducible free rotund $V$ contain an exponential point:
\[
 \exists z\in{\mathbb{C}}^n:\quad (z_1,\ldots,z_n,e^{z_1},\ldots,e^{z_n})\in V.
\]
Only existence is requested, not genericity over a coefficient field or a prescribed transcendence degree.
\par\smallskip\textit{Formulation and concept references:} \href{https://link.springer.com/article/10.1007/s00029-023-00853-y}{[S1]}.

\subsection{Short English statement}
Must every algebraic variety with the specified freedom and dimension conditions meet the graph of complex exponentiation?
\subsection{Sources}
\begin{itemize}
\item[S1]F. P. Gallinaro, Exponential sums equations and tropical geometry, Selecta Mathematica (2023).
\url{https://link.springer.com/article/10.1007/s00029-023-00853-y}.
\textit{Formulation source.}
\end{itemize}
\end{document}
